% ============================================================
%  END-OF-CHAPTER QUESTIONS
%  Sections 1-3: Member 1's questions. Sections 9-11: Member 4's.
%  Sections 4-8: written at merge (Members 2 and 3 did not supply questions).
% ============================================================
\section*{Review questions}
\addcontentsline{toc}{section}{Review questions}
\begingroup\small
\newcommand{\qsec}[1]{\par\medskip{\sffamily\bfseries\color{fxnavy}#1}\par}
\newcommand{\ans}[1]{\par{\color{fxmuted}\footnotesize\sffamily\textbf{Answer.} #1}\par\smallskip}

\qsec{Sections 1--3: Exchange rates, the forex market and the balance of payments}
\begin{enumerate}[leftmargin=1.6em, label=\arabic*., itemsep=3pt]
\item The rate moves from ₹90 to ₹95 per US dollar. Has the rupee become stronger or weaker? Explain.
\ans{Weaker. Each dollar now costs more rupees, so each rupee buys fewer dollars.}
\item Write the direct quote ``US\$1 = ₹80'' as an indirect quote.
\ans{₹1 = US\$0.0125 (1 ÷ 80).}
\item A money changer buys dollars at ₹94 and sells them at ₹96. What is the spread on US\$100?
\ans{₹200 (₹9,600 − ₹9,400).}
\item Give two sources of demand for dollars and two sources of supply of dollars in India.
\ans{Demand: importers of oil or gold, students paying foreign fees, Indian tourists abroad, Indians investing abroad. Supply: exporters, remittances, foreign tourists in India, FDI and FPI inflows.}
\item Why does a fall in the rupee against the dollar mean more rupees for Ravi, even though he is paid in dirhams?
\ans{The dirham is fixed at 3.6725 per US dollar. When the dollar costs more rupees, so does the dirham.}
\item What is the difference between FDI and FPI, and which is more likely to cause sudden moves in the rupee?
\ans{FDI buys a lasting stake in a business and stays for years. FPI buys shares and bonds and can leave in a day. FPI is more likely to cause sudden moves.}
\item Why might Selvi sign a forward contract?
\ans{To fix today the rate at which she will sell her dollars later, so that a stronger rupee does not wipe out her profit.}
\item India imports far more goods than it exports. Why is its current account deficit still small?
\ans{Large net earnings from invisibles, mainly software exports and remittances, cover most of the trade gap.}
\item How can a country pay for a current account deficit?
\ans{With a net inflow on the capital account (such as FDI, FPI or loans) or by drawing down its foreign exchange reserves.}
\end{enumerate}

\qsec{Sections 4--8: How the rate is set and managed}
\begin{enumerate}[resume, leftmargin=1.6em, label=\arabic*., itemsep=3pt]
\item At ₹80 per US dollar, buyers want US\$120 million and sellers offer US\$100 million. Is there excess demand or excess supply, and which way will the rate move?
\ans{Excess demand of US\$20 million. Buyers bid the rate up, so the rupee depreciates until the two quantities are equal.}
\item What is the difference between a depreciation and a devaluation? Give one Indian example of a devaluation.
\ans{A depreciation is a fall in a market-set rate; a devaluation is a deliberate cut in an official, fixed rate. India devalued the rupee in June 1966 (to ₹7.50 per dollar) and in July 1991.}
\item Name the three main types of exchange rate system. Which one has India used since March 1993?
\ans{Fixed (pegged), floating (flexible) and managed float. India has used a managed float since 1993.}
\item Draw a demand-and-supply diagram to show how the RBI can slow a fall in the rupee. What happens to India's reserves?
\ans{Supply of dollars shifts right as the RBI sells dollars, so the rate rises by less than it otherwise would (as in Figure~\ref{fig:intervene}). Reserves fall by the dollars sold.}
\item A pair of shoes costs ₹2,400 in India and \$40 in the US. What PPP exchange rate does this imply? If the market rate is ₹95 per dollar, is the rupee overvalued or undervalued by this measure?
\ans{₹60 per dollar (2,400 ÷ 40). At ₹95 a dollar buys more in India than PPP suggests, so the rupee is undervalued.}
\end{enumerate}

\qsec{Sections 9--11: What moves the rupee, why it matters, and crises}
\begin{enumerate}[resume, leftmargin=1.6em, label=\arabic*., itemsep=3pt]
\item The US Federal Reserve raises its interest rate sharply. Assuming nothing else changes, what is most likely to happen in the market for US dollars in India?\\
(a) The demand curve for dollars shifts left and the rupee appreciates.\\
(b) The supply curve of dollars shifts left and the rupee depreciates.\\
(c) The supply curve of dollars shifts right and the rupee appreciates.\\
(d) Nothing changes, because the RBI has not changed its own interest rate.
\ans{(b). The interest rate differential narrows, so foreign investors move funds out of India. Fewer dollars are offered at every exchange rate, and the ₹/\$ rate rises.}
\item Which of the following is an example of imported inflation?\\
(a) Vegetable prices rise after a weak monsoon.\\
(b) Diesel prices rise after the rupee falls from ₹85 to ₹90 per US dollar.\\
(c) House rents rise in Bengaluru as more IT firms open offices.\\
(d) Wages rise after a new minimum wage law.
\ans{(b). Only (b) starts with imports becoming dearer in rupees. The others are caused by domestic supply or demand conditions.}
\item In 2022, world crude oil prices rose and the US Federal Reserve raised interest rates at the same time. (i) Draw a diagram of the market for US dollars in India showing both changes. (ii) Explain why we can be sure what happened to the value of the rupee, but not to the quantity of dollars traded.
\ans{(i) Demand for dollars shifts right (refiners need more dollars for oil); supply shifts left (foreign investors take funds to the US). (ii) Both shifts push the ₹/\$ rate up, so the rupee clearly depreciates. The demand shift raises the quantity traded while the supply shift lowers it, so the net effect on quantity depends on which shift is larger.}
\item Which combination of conditions left Sri Lanka most exposed to a currency crisis in 2022?\\
(a) High reserves, little foreign debt and a booming tourism sector.\\
(b) Low reserves, heavy foreign debt repayments and a collapse in tourism earnings.\\
(c) A flexible exchange rate and fast-growing exports.\\
(d) Rising remittances and falling world oil prices.
\ans{(b). Low reserves, heavy foreign-currency debt and a sudden loss of foreign earnings are the `dry wood' described in \secref{sec:crises}.}
\end{enumerate}
\endgroup
